\section{总结与延伸}

\subsection{讲者的核心总结}

\begin{figure}[H]
\centering
\includegraphics[width=0.95\textwidth]{slides-images/slide-091.jpg}
\caption{课程总结：从预训练到 ChatGPT 的完整路径\protect\footnotemark}
\end{figure}
\footnotetext{Slides 第92页。}

Archit Sharma 在课程结尾总结了从预训练到 ChatGPT 的完整路径：

\begin{enumerate}
    \item \textbf{预训练}：在海量文本上训练 next-token prediction，获得强大的基础能力
    \item \textbf{Zero/Few-shot}：发现预训练模型具有涌现的 in-context learning 能力
    \item \textbf{SFT}：通过指令-回答对微调，使模型学会遵循指令
    \item \textbf{RLHF/DPO}：通过人类偏好数据直接优化模型行为，超越人类标注数据的质量上限
\end{enumerate}

\subsection{全课知识图谱}

\begin{center}
\begin{tikzpicture}[node distance=0.9cm, every node/.style={draw, rounded corners, minimum width=4cm, minimum height=0.7cm, font=\small}]
\node (pretrain) {预训练（Next Token Prediction）};
\node (icl) [below=of pretrain] {Zero/Few-shot + CoT};
\node (sft) [below=of icl] {指令微调（SFT）};
\node (rlhf) [below=of sft] {RLHF / DPO};
\node (chatgpt) [below=of rlhf] {ChatGPT};
\draw[->, thick] (pretrain) -- (icl);
\draw[->, thick] (icl) -- (sft);
\draw[->, thick] (sft) -- (rlhf);
\draw[->, thick] (rlhf) -- (chatgpt);
\node[draw=none, right=1cm of pretrain, text width=5.5cm, font=\scriptsize] {海量数据，学习知识、语法、推理};
\node[draw=none, right=1cm of icl, text width=5.5cm, font=\scriptsize] {无需微调，通过提示引导任务完成};
\node[draw=none, right=1cm of sft, text width=5.5cm, font=\scriptsize] {指令-回答对，多任务微调};
\node[draw=none, right=1cm of rlhf, text width=5.5cm, font=\scriptsize] {人类偏好排序，Bradley-Terry，KL 约束};
\node[draw=none, right=1cm of chatgpt, text width=5.5cm, font=\scriptsize] {对话优化，大规模部署};
\end{tikzpicture}
\end{center}

\subsection{关键 Takeaways}

\begin{importantbox}{五条核心原则}
\begin{enumerate}
    \item \textbf{预训练是基础}：几乎所有后训练的效果都依赖于强大的预训练基座。模型越大、预训练越充分，后训练的收益越大
    \item \textbf{数据质量重于数量}：无论是 SFT 还是偏好数据，高质量的少量数据可能优于低质量的大量数据（``Less Is More for Alignment''）
    \item \textbf{比较比评分容易}：人类更擅长比较两个回答的优劣，而非给单个回答打绝对分数。这一洞察是 RLHF 和 DPO 的数据基础
    \item \textbf{DPO 是 RLHF 的优雅替代}：通过巧妙的数学推导，DPO 将复杂的 RL 优化简化为简单的二分类损失，且效果相当
    \item \textbf{对齐不等于完美}：Reward hacking、冗长性、幻觉等问题提醒我们，``优化人类偏好''并不等同于``做正确的事''
\end{enumerate}
\end{importantbox}

\subsection{拓展阅读}

\begin{itemize}
    \item Radford et al., 2018. \textit{Improving Language Understanding by Generative Pre-Training} (GPT-1)
    \item Radford et al., 2019. \textit{Language Models are Unsupervised Multitask Learners} (GPT-2)
    \item Brown et al., 2020. \textit{Language Models are Few-Shot Learners} (GPT-3): \url{https://arxiv.org/abs/2005.14165}
    \item Wei et al., 2022. \textit{Chain-of-Thought Prompting Elicits Reasoning in Large Language Models}: \url{https://arxiv.org/abs/2201.11903}
    \item Kojima et al., 2022. \textit{Large Language Models are Zero-Shot Reasoners}: \url{https://arxiv.org/abs/2205.11916}
    \item Ouyang et al., 2022. \textit{Training language models to follow instructions with human feedback} (InstructGPT): \url{https://arxiv.org/abs/2203.02155}
    \item Stiennon et al., 2020. \textit{Learning to summarize from human feedback}: \url{https://arxiv.org/abs/2009.01325}
    \item Rafailov et al., 2023. \textit{Direct Preference Optimization: Your Language Model is Secretly a Reward Model} (DPO): \url{https://arxiv.org/abs/2305.18290}
    \item Zhou et al., 2023. \textit{LIMA: Less Is More for Alignment}: \url{https://arxiv.org/abs/2305.11206}
    \item Christiano et al., 2017. \textit{Deep Reinforcement Learning from Human Preferences}: \url{https://arxiv.org/abs/1706.03741}
\end{itemize}

\end{document}
